% Required packages (add to preamble if not already present):
%   \usepackage{graphicx}
%   \usepackage{subcaption}
%   \usepackage{booktabs}
%   \usepackage{siunitx}   % optional, for aligned decimal columns
%
% Adjust \graphicspath to match your thesis layout, e.g.:
%   \graphicspath{{../Results/}}
% The figure paths below assume this root.

\chapter{Evaluation}

This chapter moves from the descriptive overview of XCTest and Swift Testing established
in the preceding chapters to an empirical assessment of their practical differences.
The evaluation is structured around three areas: a qualitative analysis of the
differences, limitations, and migration experience --- encompassing the associated tooling
and test runners --- followed by a quantitative code complexity analysis across two
real-world projects, and finally a performance regression analysis aimed at detecting
any runtime overhead introduced by the transition between frameworks.

\section{Differences, Limitations, Migration and Tooling}

When comparing Swift Testing to XCTest, the differences that emerge are either
negligible from the perspective of domain-level testing, or are a consequence of
XCTest's age and the degree to which both the Swift language and software testing
practices have evolved since its introduction. Rather than introducing new concepts
or mimicking those found in XCTest, Swift Testing leans heavily on language features
and semantics, arriving with a small set of interfaces that collectively cover almost
all of the functionality found in XCTest. Swift Testing, along with its runner, is
bundled within the Xcode toolchain --- improving ergonomics of usage and removing the
need to declare it as an explicit dependency. It should be noted that the two
frameworks are not yet cross-compatible; however, the Swift Evolution proposal
ST-0021: Targeted Interoperability between Swift Testing and XCTest has been put into
motion, potentially enabling this functionality in the near future and further
lowering the migration bar.

The most prominent limitations are those concerning Swift concurrency --- in particular,
unstructured tasks created in initialisers, or functions that do not provide a
callback or return a value that could be tied to a task result. Although Swift Testing
provides the \texttt{confirmation} API, which can be combined with continuations, a
gap remains that in some cases forces the use of \texttt{Task.sleep} or
\texttt{Task.yield} in order to satisfy an assertion, producing fragile and
potentially flaky tests in the process. Similarly, certain patterns involving test
suite inheritance do not translate to Swift Testing at all, requiring a fundamentally
different approach. A less immediately visible limitation arises from the fact that,
since Xcode 26, the default actor isolation in Xcode projects has been set to
\texttt{@MainActor}, causing test suites and their contained tests to be implicitly
bound to the main actor. As a result, rather than executing in true concurrent
parallelism, tests are enqueued onto the main actor's executor and run serially with
respect to one another --- in contrast to XCTest, which spawns multiple test runners
and schedules whole test cases across them. In large codebases, this implicit
isolation could constitute a meaningful performance regression, a concern examined
further in Section~\ref{TODO}.

When it comes to migration, the official migration guide \cite{apple2024migration}
covers the majority of cases, documenting how the interfaces of the two frameworks
correspond and what is achievable with Swift Testing. The limitations noted above
notwithstanding, none present a critical obstacle: new tests can be written alongside
existing ones without additional configuration, and where a direct port proves
difficult, migration of individual test cases can be deferred to a later stage. On
Darwin platforms, where \texttt{xcodebuild} is the standard build and test tool,
both frameworks are supported out of the box --- tests from both are discovered and
executed without any additional steps. Migration to Swift Testing is therefore
relatively straightforward, best approached incrementally by introducing new test
declarations alongside the existing XCTest suite and gradually deprecating the latter,
rather than rewriting everything from scratch.

Regarding tooling, the \texttt{swift test} CLI presents several limitations that
are worth noting in the context of migration. There is currently no way to filter
test runs by tag or collection of tags directly from the command line, nor is there
support for test repetition or running until failure. The output, although significantly
improved over earlier releases, remains limited in certain respects --- most notably,
XCTest results are not surfaced in the Swift Testing output, and vice versa, meaning
that in a mixed-framework target, the results of one framework are not visible in the
other's output. This is of particular relevance in CI/CD environments, where
\cite{beller2015} demonstrated that the majority of test execution occurs, and where
result completeness is essential for reliable pipeline feedback. On Darwin platforms,
\texttt{xcodebuild} partially bridges this gap: it discovers and executes tests from
both frameworks and produces an \texttt{xcresult} bundle that captures results from
both. Community tools such as \texttt{xcbeautify} and \texttt{xcpretty} can be used
to parse and reformat this output, though they may occasionally miss individual test
case failures. The \texttt{--xunit-output} flag available via \texttt{swift test}
produces a separate \texttt{-swift-testing.xml} file alongside the XCTest output
file; however, as noted in the Methodology chapter, this output is only persisted
when the \texttt{--parallel} flag is also passed, limiting its utility under
sequential execution conditions.

\section{Code Complexity Analysis}

\subsection{Results}

% TODO: introduce the two projects briefly in one sentence each, then present tables.

\begin{table}[ht]
    \centering
    \caption{Code complexity metrics --- Project 1 (macOS Application)}
    \label{tab:complexity-project1}
    \begin{tabular}{lll}
        \hline
        \textbf{Metric}            & \textbf{XCTest} & \textbf{Swift Testing} \\
        \hline
        Lines of code              & TODO            & TODO                   \\
        Number of test functions   & TODO            & TODO                   \\
        Assertion count            & TODO            & TODO                   \\
        Boilerplate overhead (LoC) & TODO            & TODO                   \\
        \hline
    \end{tabular}
\end{table}

\begin{table}[ht]
    \centering
    \caption{Code complexity metrics --- Project 2 (SPM Framework)}
    \label{tab:complexity-project2}
    \begin{tabular}{lll}
        \hline
        \textbf{Metric}            & \textbf{XCTest} & \textbf{Swift Testing} \\
        \hline
        Lines of code              & TODO            & TODO                   \\
        Number of test functions   & TODO            & TODO                   \\
        Assertion count            & TODO            & TODO                   \\
        Boilerplate overhead (LoC) & TODO            & TODO                   \\
        \hline
    \end{tabular}
\end{table}

% TODO: prose interpreting the delta, referencing \cite{winkler2024},
% \cite{yu2019}, and the 34\% reduction reported by \cite{barbosahora2022}.

\section{Performance Regression Analysis}

\subsection{Measurement Challenges}

Establishing a reliable and comparable measurement baseline for test execution time
proved non-trivial, as neither of the two standard toolchain approaches produced
results suitable for direct comparison across frameworks.

The first approach considered was \texttt{swift test} with the
\texttt{--xunit-output} flag, which produces structured timing output in the xUnit
XML format. However, as established in the Swift Package Manager implementation,
this flag only persists output to disk when the \texttt{--parallel} flag is also
passed. Since parallel execution was explicitly excluded from this evaluation --- as
discussed in Section~\ref{TODO} --- this approach was not viable. Furthermore, when
targeting XCTest suites in isolation using \texttt{--disable-swift-testing}, no
output file is generated at all, regardless of the \texttt{--parallel} flag,
rendering this approach unsuitable for XCTest measurement as well.

The second approach involved \texttt{xcodebuild test} with the
\texttt{-resultBundlePath} flag, which captures test results in the
\texttt{xcresult} format. This worked reliably for XCTest targets; however, for
test plans containing only Swift Testing tests, the resulting \texttt{xcresult}
bundle reported test durations that were either entirely absent or fixed at a
constant value of one second, irrespective of actual execution time. This behaviour
was raised on the Swift Forums \cite{TODO-forumpost}, where Jonathan Grynspan, a
maintainer of Swift Testing at Apple, confirmed that this is a known bug in Xcode
and that it had already been reported internally.

Given that neither standard approach produced reliable per-test timing data under
the required conditions, an alternative measurement strategy was adopted, as
described in the following section.

\subsection{Custom Measurement Tool}

The empirical investigation of parallelism behaviour across both tools and their
invocation variants informed the design of a custom measurement tool,
\textit{Testbench}, developed specifically for this study. Prior to its development,
a dedicated test target was prepared containing pre-fabricated tests with known sleep
durations --- one suite sleeping for one second, one for two seconds across two tests,
and one for five seconds --- to serve as a controlled baseline for detecting whether
parallel execution was present during measurement.

Four distinct invocation strategies were evaluated. The first, \texttt{swift test}
with \texttt{--parallel} and \texttt{--num-workers}, revealed an important asymmetry
between the two frameworks. For XCTest, the number of parallel workers has a direct
and measurable effect: running three test classes with \texttt{--num-workers 3}
reduces total duration to approximately that of the longest individual class, whereas
specifying fewer workers than test classes produces proportionally longer runtimes.
For Swift Testing, the \texttt{--num-workers} flag has no observable effect --- even
with a single worker specified, tests execute in parallel, a consequence of how the
Swift Testing runner is implemented internally, independently of the worker count
passed by the host tool. Passing \texttt{--no-parallel} to \texttt{swift test}
enforces serial execution for both frameworks, producing total durations of
approximately ten seconds in both cases --- consistent with the sum of all individual
test durations.

The second strategy, \texttt{xcodebuild test}, exhibited analogous behaviour. The
\texttt{-parallel-testing-worker-count} flag has no effect on Swift Testing targets
regardless of the value supplied, while \texttt{-parallel-testing-enabled NO}
correctly enforces serial execution for both frameworks. The third strategy,
\texttt{xcodebuild build-for-testing} followed by \texttt{xcodebuild
test-without-building} with an \texttt{.xctestrun} file, behaved identically to
\texttt{xcodebuild test} with respect to parallelism: the
\texttt{-parallel-testing-worker-count} argument remained ineffective for Swift
Testing, and serial execution could again only be guaranteed via
\texttt{-parallel-testing-enabled NO}.

The fourth strategy, \texttt{xcrun xctest} invoked directly against a compiled
\texttt{.xctest} bundle, proved to be the only approach that reliably enforces serial
execution for both frameworks without additional flags, and without the timing
reliability issues identified in the preceding section. Both the XCTest and Swift
Testing bundles produced a total suite duration of approximately ten seconds ---
the sum of all individual test durations --- confirming strictly sequential execution,
as shown in Listing~\ref{lst:xcrun-xctest-output} and
Listing~\ref{lst:xcrun-swifttesting-output} respectively. The output also illustrates
the structural difference in how the two frameworks report results: XCTest emits
line-by-line output to \texttt{stderr} with explicit suite and test case boundaries,
whereas Swift Testing emits per-test and per-suite durations in its own distinct
format, also to \texttt{stderr}.

\begin{codeblock}{XCTest serial execution output of \texttt{xcrun xctest XCTestTests.xctest}}{lst:xcrun-xctest-output}{bash}
Test Suite 'All tests' started at 2026-04-11 21:07:32.681.
Test Suite 'XCTestTests.xctest' started at 2026-04-11 21:07:32.689.
Test Suite 'BarXCTests' started at 2026-04-11 21:07:32.689.
Test Case '-[XCTestTests.BarXCTests testBarOne]' started.
Test Case '-[XCTestTests.BarXCTests testBarOne]' passed (1.043 seconds).
Test Suite 'BarXCTests' passed at 2026-04-11 21:07:33.732.
     Executed 1 test, with 0 failures (0 unexpected) in 1.043 (1.043) seconds
Test Suite 'BazXCTests' started at 2026-04-11 21:07:33.732.
Test Case '-[XCTestTests.BazXCTests testBazOne]' started.
Test Case '-[XCTestTests.BazXCTests testBazOne]' passed (2.083 seconds).
Test Case '-[XCTestTests.BazXCTests testBazTwo]' started.
Test Case '-[XCTestTests.BazXCTests testBazTwo]' passed (2.020 seconds).
Test Suite 'BazXCTests' passed at 2026-04-11 21:07:37.836.
     Executed 2 tests, with 0 failures (0 unexpected) in 4.102 (4.103) seconds
Test Suite 'FooXCTests' started at 2026-04-11 21:07:37.836.
Test Case '-[XCTestTests.FooXCTests testFooOne]' started.
Test Case '-[XCTestTests.FooXCTests testFooOne]' passed (5.060 seconds).
Test Suite 'FooXCTests' passed at 2026-04-11 21:07:42.896.
     Executed 1 test, with 0 failures (0 unexpected) in 5.060 (5.060) seconds
Test Suite 'All tests' passed at 2026-04-11 21:07:42.896.
     Executed 4 tests, with 0 failures (0 unexpected) in 10.205 (10.216) seconds
\end{codeblock}

In contrast, Swift Testing emits a more compact output, reporting durations at the
per-test and per-suite level rather than interleaving start and end events, as shown
in Listing~\ref{lst:xcrun-swifttesting-output}.

\begin{codeblock}{Swift Testing serial execution output of \texttt{xcrun xctest SwiftTestingTests.xctest}}{lst:xcrun-swifttesting-output}{bash}
Test run started.
Suite BarTests started.
Test "test Bar one" passed after 1.040 seconds.
Suite BarTests passed after 1.041 seconds.
Suite BazTests started.
Test "test Baz one" passed after 2.008 seconds.
Test "test Baz two" passed after 2.029 seconds.
Suite BazTests passed after 4.039 seconds.
Suite FooTests started.
Test "test Foo one" passed after 5.001 seconds.
Suite FooTests passed after 5.001 seconds.
Test run with 4 tests in 3 suites passed after 10.082 seconds.
\end{codeblock}

Given these findings, Testbench adopts a two-phase approach. In the first phase, it
invokes \texttt{xcodebuild build-for-testing} to compile the target scheme and
produce \texttt{.xctest} bundles in a designated derived data directory. In the
second phase, it invokes \texttt{xcrun xctest} directly against each discovered
bundle, ensuring that both XCTest and Swift Testing suites are executed serially and
in isolation, without any parallelism introduced by the host tool. This guarantees
measurement conditions that are consistent and comparable across both frameworks.

Testbench is implemented as a Swift command-line tool using the \texttt{ArgumentParser}
library for its interface and the \texttt{Subprocess} library for process execution.
It exposes three subcommands: \texttt{report}, which collects and persists a machine
report including hardware, operating system, Swift, and \texttt{xcodebuild} version
information prior to any measurement run; \texttt{xcodebuild}, which orchestrates
the build-and-run cycle for a given scheme across a configurable number of
iterations, with derived data and caches cleared between runs; and \texttt{results},
which manages the stored output. Per-test and total suite execution times are
extracted from the \texttt{stderr} output of \texttt{xcrun xctest} using regular
expressions targeting the distinct output formats of each framework. Each iteration
is persisted as a structured JSON file, capturing the test plan name, framework,
total suite duration, and per-test case identifiers and durations, enabling
subsequent statistical analysis of the collected results.

\subsection{Results}
\label{sec:perf-results}

Measurements were collected on two machines: an Apple M2~Pro
(Mac14,9, 10~cores, 16~GB, Swift~6.2.3) and an Apple M5
(Mac17,2, 10~cores, 16~GB, Swift~6.2.1). Both ran macOS~Sequoia on
\texttt{arm64}. Each test suite was executed for $n = 100$ independent iterations
with derived data cleared between runs. All statistics reported below use the sample
standard deviation ($\sigma$, $\text{ddof}=1$) and a 95\% Student's
$t$-interval, which is appropriate for the sample size used.

Two projects were measured: \textit{swift-pesel}, a small SPM library whose Swift
Testing suite includes a parameterised test over 1\,000 argument cases; and
\textit{bitchat}, a macOS Bluetooth chat application whose tests involve asynchronous
delegate callbacks and structured concurrency. For \textit{bitchat}, measurements
were taken both before and after a correctness fix to the Swift Testing suite,
described in Section~\ref{sec:perf-anomaly-sleep}.

\subsubsection*{swift-pesel}

\begin{table}[ht]
    \centering
    \caption{%
        Total test run duration --- \textit{swift-pesel}
        ($n = 100$ iterations per suite; 95\% CI in parentheses).%
    }
    \label{tab:perf-pesel}
    \begin{tabular}{llrrr}
        \toprule
        \textbf{Machine} & \textbf{Suite (Framework)} &
        \textbf{Mean (ms)} & \textbf{$\sigma$ (ms)} & \textbf{95\% CI (ms)} \\
        \midrule
        M2 Pro & \texttt{PeselTests} (Swift Testing)  & 48.8 & 1.1 & [48.6,\ 49.0] \\
               & \texttt{PeselXCTests} (XCTest)        & 11.1 & 0.6 & [10.9,\ 11.2] \\
        \midrule
        M5     & \texttt{PeselTests} (Swift Testing)  & 31.5 & 0.6 & [31.3,\ 31.6] \\
               & \texttt{PeselXCTests} (XCTest)        &  7.1 & 0.8 & [\ 6.9,\ \ 7.2] \\
        \bottomrule
    \end{tabular}
\end{table}

On both machines the Swift Testing suite takes approximately $4.4\times$ longer than
its XCTest counterpart (Table~\ref{tab:perf-pesel}). This ratio, consistent across
hardware generations, is not attributable to framework overhead in isolation. The
Swift Testing suite exercises a parameterised test,
\texttt{expectValidPeselNumber(\_:)}, over 1\,000 argument cases --- each case
recorded as a distinct test event with its own result entry. The XCTest counterpart
encodes the equivalent validation as a loop within a single test method, incurring no
per-iteration tracking cost. The observed difference therefore reflects the
combination of Swift Testing's per-instance result tracking and the 1\,000-case
parameterisation, rather than any regression in single-test execution speed. The
coefficient of variation remains below 2\% for \texttt{PeselTests} and below 6\%
for \texttt{PeselXCTests}, indicating stable, low-noise measurements.

The only outlier of note is iteration~1 of \texttt{PeselXCTests} on both machines
(0.017\,s on M2~Pro, 0.015\,s on M5), which exceeds the distribution by more than
3$\sigma$ and is consistent with a cold-start effect on the first run after cache
clearance. No comparable warm-up effect is visible in subsequent iterations.

Figure~\ref{fig:pesel-boxplot} shows the marginal distributions for both suites on
both machines. The non-overlapping interquartile ranges confirm that the difference
is robust and not an artefact of measurement noise.

\begin{figure}[ht]
    \centering
    \begin{subfigure}[b]{0.48\textwidth}
        \includegraphics[width=\textwidth]{%
            m2-results/xcodebuild/swift-pesel/plots/swift-pesel_boxplot.pdf}
        \caption{Apple M2 Pro}
        \label{fig:pesel-m2-box}
    \end{subfigure}
    \hfill
    \begin{subfigure}[b]{0.48\textwidth}
        \includegraphics[width=\textwidth]{%
            m5-results/xcodebuild/swift-pesel/plots/swift-pesel_boxplot.pdf}
        \caption{Apple M5}
        \label{fig:pesel-m5-box}
    \end{subfigure}
    \caption{%
        Total test run duration for \textit{swift-pesel} on both machines
        (box: Q1--Q3; whiskers: $1.5\times\mathrm{IQR}$; $\blacklozenge$: mean;
        $n=100$ per suite).
    }
    \label{fig:pesel-boxplot}
\end{figure}

\subsubsection*{bitchat --- pre-fix}

\begin{table}[ht]
    \centering
    \caption{%
        Total test run duration --- \textit{bitchat}, before correctness fix
        ($n = 100$ iterations per suite; 95\% CI in parentheses).%
    }
    \label{tab:perf-bitchat-pre}
    \begin{tabular}{llrrr}
        \toprule
        \textbf{Machine} & \textbf{Suite (Framework)} &
        \textbf{Mean (s)} & \textbf{$\sigma$ (s)} & \textbf{95\% CI (s)} \\
        \midrule
        M2 Pro & \texttt{bitchatTests} (Swift Testing)  & 8.330 & 0.117 & [8.307,\ 8.353] \\
               & \texttt{bitchatXCTests} (XCTest)        & 5.023 & 0.051 & [5.013,\ 5.034] \\
        \midrule
        M5     & \texttt{bitchatTests} (Swift Testing)  & 7.946 & 0.026 & [7.941,\ 7.951] \\
               & \texttt{bitchatXCTests} (XCTest)        & 4.775 & 0.026 & [4.770,\ 4.781] \\
        \bottomrule
    \end{tabular}
\end{table}

Prior to the fix, the Swift Testing suite ran 3.31\,s slower than XCTest on the M2~Pro
($+65.9\%$) and 3.17\,s slower on the M5 ($+66.3\%$), as shown in
Table~\ref{tab:perf-bitchat-pre} and Figure~\ref{fig:bitchat-pre-box}. The
distributions are well-separated and non-overlapping, ruling out any plausible
measurement artefact. The gap magnitude is also consistent across both hardware
platforms, which strengthens the case that the cause is specific to the suite's
implementation rather than to system-level scheduling noise.

\begin{figure}[ht]
    \centering
    \begin{subfigure}[b]{0.48\textwidth}
        \includegraphics[width=\textwidth]{%
            m2-pro-results-pre-swift-testing-refactor/xcodebuild/bitchat/plots/bitchat_boxplot.pdf}
        \caption{Apple M2 Pro --- before fix}
        \label{fig:bitchat-pre-m2-box}
    \end{subfigure}
    \hfill
    \begin{subfigure}[b]{0.48\textwidth}
        \includegraphics[width=\textwidth]{%
            m5-pro-results-pre-swift-testing-refactor/xcodebuild/bitchat/plots/bitchat_boxplot.pdf}
        \caption{Apple M5 --- before fix}
        \label{fig:bitchat-pre-m5-box}
    \end{subfigure}
    \caption{%
        Total test run duration for \textit{bitchat} before the correctness fix,
        on both machines. The Swift Testing suite is approximately 3.2--3.3\,s
        slower than XCTest, with non-overlapping interquartile ranges on both
        platforms ($n = 100$ per suite).
    }
    \label{fig:bitchat-pre-box}
\end{figure}

\subsubsection*{bitchat --- root cause of the pre-fix gap}
\label{sec:perf-anomaly-sleep}

Root-cause analysis of the pre-fix gap identified a systematic misuse of Swift
Testing's \texttt{confirmation} API in \texttt{BLEServiceTests.swift}. Seven
asynchronous tests --- covering \texttt{sendMessage}, \texttt{sendPrivateMessage},
\texttt{sendMessageWithMentions}, \texttt{simulateIncomingMessage},
\texttt{simulateIncomingPacket}, \texttt{messageDeliveryHandler}, and
\texttt{packetDeliveryHandler} --- each contained an unconditional
\texttt{sleep(0.5)} call inside their \texttt{confirmation} body, totalling 3.5\,s
of dead wait per run:

\begin{verbatim}
@Test func sendPublicMessage() async throws {
    try await confirmation { receivedPublicMessage in
        ...
        service.sendMessage("Hello, world!")
        try await sleep(0.5)   // ← 500 ms regardless of when callback fires
    }
}
\end{verbatim}

The XCTest counterpart used \texttt{fulfillment(of:timeout:)}, which returns as soon
as the expectation is fulfilled --- in practice, sub-millisecond once the main queue
processes the pending \texttt{DispatchQueue.main.async} block:

\begin{verbatim}
func testSendPublicMessage() async {
    let exp = expectation(description: "...")
    ...
    service.sendMessage("Hello, world!")
    await fulfillment(of: [exp], timeout: 1.0)  // ← returns on fulfillment
}
\end{verbatim}

The critical semantic difference is that \texttt{confirmation} has no built-in
polling: it executes the body closure synchronously to completion and then asserts
the callback count. Any \texttt{sleep} inside it constitutes unconditional dead time,
whereas \texttt{fulfillment(of:timeout:)} drives the run loop and returns
immediately on fulfilment. Of the seven affected tests, only three actually required
any sleep --- those where the delegate callback is dispatched via
\texttt{DispatchQueue.main.async} in \texttt{MockBLEService.swift} (lines~75 and~104).
The remaining four observe synchronous callback paths
(\texttt{simulateIncomingMessage}, \texttt{simulateIncomingPacket},
\texttt{packetDeliveryHandler}) and require no sleep at all, since the confirmation
fires before the body closure returns.

The fix reduced the three necessary sleeps from 500\,ms to 10\,ms --- sufficient for
a single main-queue dispatch to drain --- and removed the four unnecessary ones
entirely, eliminating the $7 \times 500\,\mathrm{ms} = 3.5\,\mathrm{s}$ overhead.

\subsubsection*{bitchat --- post-fix}

\begin{table}[ht]
    \centering
    \caption{%
        Total test run duration --- \textit{bitchat}, after correctness fix
        ($n = 100$ iterations per suite; 95\% CI in parentheses).%
    }
    \label{tab:perf-bitchat-post}
    \begin{tabular}{llrrr}
        \toprule
        \textbf{Machine} & \textbf{Suite (Framework)} &
        \textbf{Mean (s)} & \textbf{$\sigma$ (s)} & \textbf{95\% CI (s)} \\
        \midrule
        M2 Pro & \texttt{bitchatTests} (Swift Testing)  & 4.556 & 0.036 & [4.548,\ 4.563] \\
               & \texttt{bitchatXCTests} (XCTest)        & 4.995 & 0.033 & [4.988,\ 5.001] \\
        \midrule
        M5     & \texttt{bitchatTests} (Swift Testing)  & 4.351 & 0.040 & [4.343,\ 4.359] \\
               & \texttt{bitchatXCTests} (XCTest)        & 4.783 & 0.037 & [4.776,\ 4.791] \\
        \bottomrule
    \end{tabular}
\end{table}

Following the fix, the Swift Testing suite completes in 4.556\,s on the M2~Pro and
4.351\,s on the M5 --- in both cases \textit{faster} than its XCTest counterpart
by approximately 0.44\,s ($-8.8\%$) and 0.43\,s ($-9.0\%$) respectively
(Table~\ref{tab:perf-bitchat-post}, Figure~\ref{fig:bitchat-post-box}). The
XCTest baseline remained unchanged between the two measurement rounds within the
bounds of normal variation ($\Delta = -0.028\,\mathrm{s}$ on M2~Pro,
$\Delta = +0.008\,\mathrm{s}$ on M5), confirming that the change was isolated
to the Swift Testing suite.

The remaining difference in favour of Swift Testing may be attributable to
differences in per-test isolation mechanics --- Swift Testing uses value-typed
\texttt{struct} suites with per-test initialisation, whereas XCTest allocates a
fresh \texttt{class} instance and calls \texttt{setUp}/\texttt{tearDown} for
each test, which can carry higher allocation and Objective-C messaging overhead.
The precise contribution of this factor is outside the scope of this study.

\begin{figure}[ht]
    \centering
    \begin{subfigure}[b]{0.48\textwidth}
        \includegraphics[width=\textwidth]{%
            m2-results/xcodebuild/bitchat/plots/bitchat_boxplot.pdf}
        \caption{Apple M2 Pro --- after fix}
        \label{fig:bitchat-post-m2-box}
    \end{subfigure}
    \hfill
    \begin{subfigure}[b]{0.48\textwidth}
        \includegraphics[width=\textwidth]{%
            m5-results/xcodebuild/bitchat/plots/bitchat_boxplot.pdf}
        \caption{Apple M5 --- after fix}
        \label{fig:bitchat-post-m5-box}
    \end{subfigure}
    \caption{%
        Total test run duration for \textit{bitchat} after the correctness fix.
        The Swift Testing suite is now marginally but consistently faster than
        XCTest on both platforms ($n = 100$ per suite).
    }
    \label{fig:bitchat-post-box}
\end{figure}

Figure~\ref{fig:bitchat-ts-m2} shows the per-iteration time series for the
post-fix \textit{bitchat} measurements on the M2~Pro. Both suites are stable
across all 100 iterations with no visible trend, drift, or warm-up effect,
supporting the validity of the mean as a summary statistic.

\begin{figure}[ht]
    \centering
    \includegraphics[width=\textwidth]{%
        m2-results/xcodebuild/bitchat/plots/bitchat_timeseries.pdf}
    \caption{%
        Per-iteration total run duration for \textit{bitchat} (post-fix,
        Apple M2~Pro). Dashed lines indicate the per-suite mean. Both
        suites exhibit stable, low-variance behaviour across all 100 iterations.
    }
    \label{fig:bitchat-ts-m2}
\end{figure}

\subsubsection*{Summary}

The results confirm that, in the absence of implementation defects, Swift Testing
introduces no performance regression relative to XCTest. For the small-suite
\textit{swift-pesel} project, the $4.4\times$ overhead of the Swift Testing suite
is attributable to the framework's per-instance result tracking over a
1\,000-case parameterised test, not to any regression in single-test overhead. For
the larger \textit{bitchat} project, a superficially large gap --- $+66\%$ on both
machines --- was caused entirely by seven unconditional \texttt{sleep(0.5)} calls
inside \texttt{confirmation} bodies in the Swift Testing suite, a misuse of the API
with no counterpart in the XCTest suite. After correcting this pattern, Swift Testing
ran approximately 9\% \textit{faster} than XCTest on both platforms. The finding
corroborates the key distinction between \texttt{confirmation} and
\texttt{fulfillment(of:timeout:)}: the latter exits as soon as the expectation is
fulfilled, while the former always runs the body to completion, making any sleep
inside it unconditional dead time. This difference is easy to overlook during
migration and represents a meaningful performance pitfall when porting
asynchronous tests.

\subsection{Anomalies}

Two anomalies were identified during the course of the evaluation. The first
concerns parameterised tests executed directly from the Xcode IDE: when a
parameterised test function is supplied with a sufficiently large argument
collection, Xcode tends to become unresponsive for an extended period while
awaiting the test result. This behaviour was not observed when invoking the same
test via \texttt{swift test} or \texttt{xcodebuild}, isolating the issue to the
IDE layer rather than the framework or its runner. The issue has been filed with
Apple as FB13986710 and remains unresolved at the time of writing.

The second concerns the \texttt{@MainActor} isolation behaviour introduced as the
default in Xcode 26 projects, discussed in Section~\ref{TODO}. While its impact
was not directly measurable within the scope of the two projects examined here,
the architectural implication --- that tests bound to the main actor executor run
serially rather than in parallel --- warrants attention in the context of larger
codebases, where the performance gap relative to XCTest's multi-runner scheduling
model may become significant.
